% Generated draft: requires mathematical and source review.
\documentclass[11pt]{article}
\usepackage[margin=25mm]{geometry}
\usepackage{fontspec}
\setmainfont{Latin Modern Roman}
\usepackage{amsmath,amssymb,mathtools,unicode-math}
\setmathfont{Latin Modern Math}
\usepackage{enumitem,needspace,xurl,hyperref}
\setlength{\parindent}{0pt}
\setlength{\parskip}{5pt}
\setlength{\emergencystretch}{3em}
\newcommand{\N}{\mathbb{N}}
\newcommand{\Z}{\mathbb{Z}}
\newcommand{\Q}{\mathbb{Q}}
\newcommand{\R}{\mathbb{R}}
\newcommand{\C}{\mathbb{C}}
\newcommand{\F}{\mathbb{F}}
\newcommand{\A}{\mathbb{A}}
\newcommand{\PP}{\mathbb{P}}
\newcommand{\E}{\mathbb{E}}
\newcommand{\eps}{\varepsilon}
\newcommand{\dd}{\,\mathrm{d}}
\begin{document}
\section*{UnsolvedMath ID: 11000043}
\textbf{Source problem code:} catalogue entry
\subsection{Definitions and mathematical statement}
For \(g\geq2\), let \(\mathcal M_g=\operatorname{Teich}_g/\operatorname{Mod}_g\) be the moduli orbifold of closed oriented genus-\(g\) surfaces, in the usual convention. As the source's footnote specifies, an orbifold Riemannian metric is studied via a \(\operatorname{Mod}_g\)-invariant Riemannian metric on \(\operatorname{Teich}_g\).

\emph{Conjecture 4.4 (nonpositive curvature).} There is no complete finite-volume Riemannian metric \(h\) on \(\mathcal M_g\) for which sectional curvatures are nonpositive and uniformly bounded below. Explicitly, for every \(g\geq2\), no such \(h\) and finite constant \(C\geq0\) satisfy
\[
 \operatorname{Vol}_h(\mathcal M_g)<\infty,
 \qquad -C\leq K_h(\Pi)\leq0
\]
for every tangent \(2\)-plane \(\Pi\), together with metric completeness. The lower bound is allowed to depend on the metric and genus but must hold uniformly throughout the space. The statement does not require a negative upper bound separated from zero.

The source suggests stronger possible versions: relax the finite-volume hypothesis, remove the uniform lower curvature bound, or extend the nonexistence statement to every finite cover and perhaps certain infinite covers. It also proposes an analogue for CAT(0) metrics using a generalized finite-volume notion from [BrF]. Those are suggestions, not included hypotheses or established extensions. The orbifold/cover formulation and the precise CAT(0) finite-volume convention must be specified.

The surrounding discussion points to relative curvature notions [MM1, MM2], weaker holomorphic or Ricci curvature and highly singular metrics [LSY], and Kobayashi or K\"ahler hyperbolicity [Mc2], and identifies cusp topology as an obstacle to standard metric models. The next section, ``Scalar curvature and Q-rank,'' sets \(d(S)=3g-3+n\) for a genus-\(g\), \(n\)-punctured surface and describes it as Teichm\"uller complex dimension and the pants-decomposition curve count. Those descriptions require the usual stable-surface conventions, absent from that excerpt. It mentions work of Shmuel Weinberger and the author [FW2], but the theorem itself is not supplied and is not reconstructed here.
\subsection{Short English statement}
For every genus \(g\geq2\), the conjecture rules out a complete finite-volume moduli-orbifold metric with all sectional curvatures between a finite uniform lower bound and zero. The optional stronger versions and later scalar-curvature discussion are kept separate.
\subsection{Sources}
\begin{itemize}
\item Local UnsolvedMath registry, permanent ID 11000043, source code catalogue entry, Conjecture 4.4, ``Nonpositive curvature,'' PDF page 38. The text identifies the surrounding author as B. Farb; no proposer field is supplied.
\item Farb, \emph{Problems on Mapping Class Groups and Related Topics}. Supplied URL: \url{https://www.math.uchicago.edu/~farb/papers/mcgbook.pdf}; recorded access 2026-07-29; extraction pdf-font-structure. The edited volume contains 329 explicitly typeset headings across contributed chapters, with repeated local numbering disambiguated by page. Footnote 5 explains the invariant-Teichm\"uller-metric convention.
\item Internal references [BrF], [MM1, MM2], [LSY], [Mc2], and [FW2] occur in the contextual discussion, with Shmuel Weinberger credited in the last. Full bibliography and the following theorem are absent.
\item Supplied status open and OPEN-TRIAGE. The summary offers only a general observation that metric-curvature questions are subtle and mentions unspecified mixing/non-NPC results, even saying ``Without statement, mark open/pending.'' It provides no precise assessment of the actual supplied conjecture. Current primary-literature status and redistribution terms remain NEEDS\_REVIEW.
\end{itemize}
\end{document}
